% Original PDF: source/普通高考/2015/2015北京理.pdf, page 1, question 3.
% pdfimages -list returned no embedded bitmap; the program flowchart is vector/text artwork.
% Grid evidence: tmp/redraw_flowchart_2015_beijing_science/source/page1_grid100.jpg;
% q3_block_grid50.jpg; comparison source is q3_block.png from the no-grid page render.
% Elements: rounded 开始/结束 terminals; assignment boxes
% x=1,y=1,k=0; s=x-y,t=x+y; x=s,y=t; k=k+1; decision k>=3;
% output parallelogram 输出 (x,y); solid directed connectors and 是/否 labels.
% Node -> directed edge list: 开始->(x=1,y=1,k=0)->(s=x-y,t=x+y)
% ->(x=s,y=t)->(k=k+1)->(k>=3); (k>=3)-是->输出 (x,y)->结束;
% (k>=3)-否->the right return branch->the connector entering s=x-y,t=x+y.
% solid segments: every node border, trunk arrows, right return branch, and arrowheads;
% dashed segments: none. Labels are centered in nodes, locally zero-indented,
% with compact white margins; 是 is beside the output branch and 否 beside the return branch.

\begin{tikzpicture}[
  >=stealth,
  line width=0.55pt,
  line cap=round,
  line join=round,
  every node/.style={anchor=center,align=center,
    execute at begin node={\setlength{\parindent}{0pt}}},
  flowbox/.style={draw,line width=0.55pt,minimum height=0.70cm,
    inner xsep=4pt,inner ysep=2pt},
  io/.style={flowbox,shape=trapezium,
    trapezium left angle=72,trapezium right angle=108,
    minimum height=0.78cm},
  flowarrow/.style={->,line width=0.55pt}
]
  % Main vertical flow; the loop returns to the connector above the s,t box.
  \node[flowbox,rounded corners=0.18cm,minimum width=1.35cm]
    (start) at (0,9.20) {开始};
  \node[flowbox,minimum width=3.65cm]
    (init) at (0,7.95) {$x=1, y=1, k=0$};
  \node[flowbox,minimum width=3.72cm]
    (st) at (0,6.62) {$s=x-y, t=x+y$};
  \node[flowbox,minimum width=2.65cm]
    (assign) at (0,5.25) {$x=s, y=t$};
  \node[flowbox,minimum width=2.35cm]
    (increment) at (0,3.90) {$k=k+1$};
  \node[flowbox,shape=diamond,aspect=2.75,minimum width=3.95cm,
    minimum height=1.02cm]
    (test) at (0,2.50) {$k\geq 3$};
  \node[io,minimum width=2.75cm]
    (output) at (0,1.15) {输出 $(x,y)$};
  \node[flowbox,rounded corners=0.18cm,minimum width=1.35cm]
    (finish) at (0,-0.08) {结束};

  \draw[flowarrow] (start.south) -- (init.north);
  \draw[flowarrow] (st.south) -- (assign.north);
  \draw[flowarrow] (assign.south) -- (increment.north);
  \draw[flowarrow] (increment.south) -- (test.north);
  \draw[flowarrow] (test.south) -- (output.north);
  % The trapezium's south anchor is offset by its slanted side; keep this
  % explicit central connector to avoid a false diagonal edge.
  \draw[flowarrow] (0,0.77) -- (finish.north);

  % The no branch goes right, up, then left into the s,t connector.
  \draw[flowarrow] (test.east) -- (3.30,2.50) -- (3.30,7.16) -- (0,7.16);
  % Main connector continues from the loop junction into the s,t box.
  \draw[flowarrow] (init.south) -- (0,7.16) -- (st.north);

  \node[anchor=west,inner sep=0pt] at (0.48,1.90) {是};
  \node[anchor=west,inner sep=0pt] at (2.42,2.92) {否};

  \path[use as bounding box] (-2.15,9.55) rectangle (4.05,-0.42);
\end{tikzpicture}
